固定负荷下的储能调度可写成连续能源流优化，并通过购售电价格与荷电状态约束刻画跨时价值\cite{zhang_dc_bess_2025}。对于线性成本模型中可能出现的同时充放电现象，可通过次级目标消除无意义循环\cite{wang_exact_relax_2024}。因此采用“先最小成本、再最小吞吐量”的两层字典序求解。

\begin{algorithm}[H]
\SetAlgoLined
\setstretch{0.86}
\SetAlFnt{\small}
\SetArgSty{textnormal}
\SetKwInput{KwIn}{输入}
\SetKwInput{KwInit}{初始化}
\SetKwInput{KwOut}{输出}
\SetKwFor{ForEach}{对每个}{执行}{结束}
\caption{固定负荷下的两层字典序储能经济调度}
\label{alg:q3-two-stage}
\KwIn{设施负荷 $L_{rt}$，可用新能源 $R_{rt}$，购售电价格、碳强度，储能与电网参数}
\KwInit{$\varepsilon\leftarrow10^{-4}$ CNY；设置附件原始方案、无储能优化方案和 BESS-ED 方案}

\ForEach{区域 $r\in\{A,B,C,D,E,F\}$}{
  读取 $L_{rt}$、$R_{rt}$、价格及储能参数，构造 $u,qR,qG,d,gL,s,w,E$ 的变量索引\;
  建立新能源平衡、负荷平衡和 SOC 递推等式，并加入充放电功率、购售电上限及终端 SOC 约束\;

  \textbf{第一层：最小化净运行成本}\;
  构造
  $\mathrm{Cost}_r=\sum_t[c^{buy}_{rt}(gL_{rt}+qG_{rt})-c^{sell}_{rt}s_{rt}]\Delta t$\;
  求解区域线性规划；若不可行则返回约束信息，否则记录最优值 $\mathrm{Cost}_r^*$\;

  \textbf{第二层：在成本最优集合内最小化储能吞吐量}\;
  增加 $\mathrm{Cost}_r\le\mathrm{Cost}_r^*+\varepsilon$\;
  最小化 $\mathrm{Throughput}_r=\sum_t(qR_{rt}+qG_{rt}+d_{rt})\Delta t$\;

  重构 $G^{buy}_{rt}=gL_{rt}+qG_{rt}$、$P^{net}_{rt}=G^{buy}_{rt}-s_{rt}$ 与逐时 SOC\;
  计算区域成本、碳排放、新能源利用率、峰值净购电、净电网波动、平均绝对爬坡及储能吞吐量\;
  检查新能源/负荷平衡、SOC 递推、终端 SOC、充放电功率、购售电边界与同时充放电小时数\;
}

汇总六区域结果，并以“储能协同方案减无储能优化方案”的差值计算储能增量价值\;
\KwOut{六区域逐时储能与购售电方案、系统评价指标及约束检验结果}
\end{algorithm}
